\section{Main Objectives}

\begin{frame}
\frametitle{Quantum theory of raidation by nonstationary systems 
	with application to high-order hamronic generation
        (Phys. Rev. A 101, 013410 (2020))}
	\framesubtitle{Overview of the problem}
\footnotesize
\begin{textblock*}{400pt}(30pt,60pt)
\begin{itemize}
\item We are ultimately interested in solving 
the following 1D problem with a ZRP
\begin{equation}
	i\frac{\partial}{\partial t}\phi_\omega(x,t)
	   = \left( \hat{H}_a + \hat{V}_{\rm ext}(t)\right) 
	   \phi_\omega (x,t)
 	   + e^{i\omega t} {\hat{p}}\,
	   \psi(x,t), \label{eqn:main_1d}
\end{equation}
%\begin{equation}
%	i\frac{\partial}{\partial t}\phi_\omega(x,t)
%	   = \left( \hat{H}_a + \hat{V}_{\rm ext}(t)\right) 
%	   \phi_\omega (x,t)
% 	   + \alpha A_\omega e^{i\omega t} {\hat{p}}\,
%	   \psi(x,t), \label{eqn:main_1d}
%\end{equation}

where

\begin{equation}
  \hat{H}_a = -\frac{1}{2}\frac{d^2}{dx^2} - \varkappa\delta(x), 
\hspace{20pt}  \hat{V}_{\rm ext} = F(t)x,
\end{equation}
\begin{equation}
  \phi_\omega (x,0) = \frac{-i\varkappa^{3/2}}{\omega}
{\rm sgn}(x)\left(e^{-\varkappa|x|} - e^{-\varkappa(\omega)|x|}\right), 
\hspace{20pt}  \psi(x,0) = \varkappa^{1/2}e^{-\varkappa|x|}, 
\label{eqn:bound_state0}
\end{equation}

with

\begin{equation}
  \hat{p} = -id/dx,\hspace{10pt} \mbox{ and } 
  \hspace{10pt} \varkappa(\omega) = \sqrt{\varkappa^2 + 2\omega}.
\end{equation}
%\begin{equation}
%  E_k = k^2/2,\hspace{20pt}
%  \psi_k^{(\pm)}(x) = e^{ikx} + R(\pm|k|)e^{\pm i|kx|}, \hspace{20pt}
%   R(k) = \frac{i\varkappa}{k-i \varkappa}. \label{eqn:inout_states}
%\end{equation}
%\item The parameters of the laser we will be using 
%are given by the relations
%\begin{equation}
%  F(t) = F_0 e^{-(2t^\prime/\tau)^2}\cos(\omega_0 t^\prime),\hspace{20pt}
%  0 \leqslant t^\prime \leqslant \tau
%\end{equation}
%\begin{equation}
%  \tau  = 2\pi n_{oc}/\omega,  \hspace{20pt}
%  F_0 = .1\,{\rm a.u.}, \hspace{20pt}
%  n_{oc} = 5, \hspace{20pt}
%  \omega_0 = 0.114\,{\rm a.u.}
%\end{equation}
\end{itemize}

\end{textblock*}

\begin{textblock*}{250pt}(240pt,80pt)
%\includegraphics[scale=.25]{dip_app_limits.png}{\footnote{\Tiny
%H.R. Reiss, Phys Rev A 63 013409 (2000)}}
\end{textblock*}

\end{frame}





\begin{frame}
\frametitle{Quantum theory of raidation by nonstationary systems 
	with application to high-order hamronic generation
        (Phys. Rev. A 101, 013410 (2020))}
        \framesubtitle{Overview of the problem}
\footnotesize
\begin{textblock*}{400pt}(30pt,60pt)
\begin{itemize}
\item To this end, we need to resolve $\psi(x,t)$ 
which is given by the regular TDSE, i.e 
\begin{equation}
	i\frac{\partial}{\partial t}\psi(x,t)
	   = \left( \hat{H}_a + \hat{V}_{\rm ext}(t)\right) 
	   \psi(x,t),   \hspace{20pt}\psi(x,0)=\psi_0(x),  \label{eqn:step_1d}
\end{equation}
\begin{equation}
  \hat{H}_a = -\frac{1}{2}\frac{d^2}{dx^2} - \varkappa\delta(x), \hspace{20pt}
  \hat{V}_{\rm ext} = F(t)x,
\end{equation}
\begin{equation}
  E_0 = -\varkappa^2/2,\hspace{20pt}
  \psi_{0}(x) = \varkappa^{1/2}e^{-\varkappa|x|}, \hspace{20pt}
  \varkappa = 1       \label{eqn:bound_state}
\end{equation}
\begin{equation}
  E_k = k^2/2,\hspace{20pt}
  \psi_k^{(\pm)}(x) = e^{ikx} + R(\pm|k|)e^{\pm i|kx|}, \hspace{20pt}
   R(k) = \frac{i\varkappa}{k-i \varkappa}. \label{eqn:inout_states}
\end{equation}
\vspace{10pt}
\item The parameters of the laser we will be using 
are given by the relations
\begin{equation}
  F(t) = F_0 e^{-(2t^\prime/\tau)^2}\cos(\omega_0 t^\prime),
\hspace{20pt}  0 \leqslant t \leqslant T,
\hspace{20pt}  t^{\prime} = t - T/2,
\end{equation}
\begin{equation}
  \tau  = 2\pi n_{oc}/\omega_0,  \hspace{20pt}
  n_{oc} = 5. \hspace{20pt}
\end{equation}
\end{itemize}

\end{textblock*}

\begin{textblock*}{250pt}(240pt,80pt)
%\includegraphics[scale=.25]{dip_app_limits.png}{\footnote{\Tiny
%H.R. Reiss, Phys Rev A 63 013409 (2000)}}
\end{textblock*}

\end{frame}




\begin{frame}
\frametitle{Quantum theory of raidation by nonstationary systems 
	with application to high-order hamronic generation
        (Phys. Rev. A 101, 013410 (2020))}
	\framesubtitle{Survival probability and PEMD}

\begin{textblock*}{400pt}(30pt,60pt)
\footnotesize
\begin{itemize}
\item The coefficients of the projection of $\bar{\Psi}$ are then, at $T$,
and by definition, expressed as
\begin{equation}
a_0 = e^{iE_0 T}\int \psi_0(x) \psi(x,T)dx, \hspace{20pt}
a_k = e^{iE_k T}\int \psi_k^{\Tiny (-)*}(x) \psi(x,T)dx. \hspace{20pt}
\label{eqn:a0k}
\end{equation}

%with 
%$\bar{\Psi}_{k_j} = a_{k_j} 
%\hspace{10pt} 0\leqslant j \leqslant m$, \hspace{30pt}
%$\bar{\Psi} = \left( \bar{\Psi}_0, \bar{\Psi}_{k_1}, \cdots, \bar{\Psi}_{k_j}
%\cdots, {\Psi}_{k_m}\right)$.


\item The survival probability is then expressed, 
at the end of the interaction, by
\begin{align}
   P_0         & = \left| a_0 \right|^2. \label{eqn:surv}  
\end{align}

\item The ionization probability can then be found many different 
ways that should be consistent with each others
\begin{align}
   P_{\rm ion} & = \int P(k)\frac{dk}{2\pi}, 
 \hspace{30pt} \mbox{with}\hspace{30pt} P(k) = \left| a_k \right|^2. 
\label{eqn:pion1} \\
               & = 1-P_0.
\end{align}

\item  $P(k)$ is the probability for the electron 
being ionized with the energy $E_k$.

\end{itemize}
\end{textblock*}


\end{frame}



\begin{frame}
\frametitle{Quantum theory of raidation by nonstationary systems 
	with application to high-order hamronic generation
        (Phys. Rev. A 101, 013410 (2020))}
	\framesubtitle{Quantum spectrum of HHG photons (semi-classical approx.) - 1}

\begin{textblock*}{400pt}(30pt,60pt)
\footnotesize
\begin{itemize}
\item In the classical approx., the spectrum is obtained from the following 
(Eq.(80))
\begin{equation}
Q^{(c)}_\nu (\omega) = 
 \alpha^2 \omega^2
    \left| \boldsymbol{\rm A}_{\nu}\boldsymbol{\rm d}(\omega) \right|^2,
\label{eqn:Qc}
\end{equation}
where for an arbitrary operator $\hat{f}$, we have   
\begin{equation}
f(t) = \int \psi^{*}(x,t)\hat{f}\,\psi(x,t)dx
\label{eqn:mean_value}
\end{equation}
and the Fourier transform is obtained using 
\begin{equation}
f(\omega) = \int f(t)e^{i\omega t}dt.
\label{eqn:fourier}
\end{equation}

\item Also, we consider the relations
\begin{align}
% \boldsymbol{\rm A}_\nu 
%   = \sqrt{\frac{2\pi c^2}{\omega L^3}}\boldsymbol{\rm e}_\nu,
%   \hspace{20pt}\alpha = 1/c.
 \hat{p} = -i d/dx, \hspace{20pt} \hat{d} = -x
\end{align}
\end{itemize}
\end{textblock*}


\end{frame}






\begin{frame}
\frametitle{Quantum theory of raidation by nonstationary systems 
	with application to high-order hamronic generation
        (Phys. Rev. A 101, 013410 (2020))}
	\framesubtitle{Quantum spectrum of HHG photons (semi-classical approx.) - 2}

\begin{textblock*}{400pt}(30pt,60pt)
\footnotesize
\begin{itemize}
\item Let us apply the following simplifications
\begin{align}
    \nu \longrightarrow \omega, \hspace{20pt}
    \hspace{20pt} \boldsymbol{\rm A}_\nu \longrightarrow A_\omega, 
    \hspace{20pt}\mbox{ and } \hspace{20pt} \alpha A_\omega = 1.
\end{align}
\item So the observable we are considering here is 
\begin{align}
    Q^{(c)}(\omega) = \omega^2 \left| d(\omega)  \right|^2.
\end{align}
\item Additionally let us consider Eq. (104)
\begin{align}
    Q^{(c)}(\omega) = \left| p(\omega)  \right|^2.
\end{align}
and, by virtue of $\boldsymbol{p}(t)= -\dot{\boldsymbol{d}}(t)$, which reads
\begin{align}
    Q^{(c)}(\omega) = \left| -\dot{d}(\omega)  \right|^2. 
\label{eqn:dipole_mom_diff}
\end{align}

\end{itemize}
\end{textblock*}


\end{frame}


\begin{frame}
\frametitle{Quantum theory of raidation by nonstationary systems 
	with application to high-order hamronic generation
        (Phys. Rev. A 101, 013410 (2020))}
	\framesubtitle{Length, velocity and acceleration gauges}

\begin{textblock*}{400pt}(30pt,60pt)
\footnotesize
\begin{itemize}
\item (\ref{eqn:mean_value}) applied to the operator $\hat{d}$ gives
its expression in the length gauge
\begin{equation}
d(t) = \langle \hat{d\,} \rangle = - \left\langle\hat{x}\right\rangle
  = - \int \psi^{*}(x,t)\hat{x}\,\psi(x,t)dx
\end{equation}

\item $\dot{d}(t)$ is given by the formula
\begin{equation}
\dot{d}(t) = \frac{d}{dt}d(t)
\end{equation}

\item $p(t)$ is given by the formula
\begin{equation}
p(t) = \langle \hat{p\,} \rangle  
  = - i\int \psi^{*}(x,t)\frac{d}{dx}\psi(x,t)dx
\end{equation}

\item Using the Ehrenfest theorem for the operator $x$, we get the relations 
\begin{align}
\frac{d}{dt}\langle \hat{x\,} \rangle 
        & = \frac{\langle \hat{p\,} \rangle}{m}, \\
\frac{d}{dt} \langle \hat{p\,}\rangle 
        & = \left\langle -\frac{\partial V_{\rm ext}}{\partial x} \right\rangle
\end{align}

and some others (see Bransden and Joachain p.531-532) 



\end{itemize}
\end{textblock*}


\end{frame}




\begin{frame}
\frametitle{Quantum theory of raidation by nonstationary systems 
	with application to high-order hamronic generation
        (Phys. Rev. A 101, 013410 (2020))}
	\framesubtitle{Additional formula from Ehrenfest theorem}

\begin{textblock*}{400pt}(30pt,60pt)
\footnotesize
\begin{itemize}
\item  Using the integration by parts formula, we get
\begin{equation}
   p(\omega) = \int_{-T/2}^{T/2} e^{i\omega t} \frac{dx(t)}{dt}dt,\hspace{30pt}
   u(t) = x(t),\;\;\; v(t) = e^{i\omega t},
\end{equation}
\noindent thus
\begin{align}
   p(\omega) & = \int_{-T/2}^{T/2} \frac{d}{dt}\Big(e^{i\omega t}x(t)\Big)dt
    -i\omega \int_{-T/2}^{T/2} e^{i\omega t}x(t)dt, \\
             & = \Bigg[ e^{i\omega t} x(t)\Bigg]_{-T/2}^{T/2}
               -i\omega \int_{-T/2}^{T/2} e^{i\omega t}x(t)dt. 
\label{eqn:additional}
\end{align}

So, we get an additional formula for evaluating 
\begin{equation}
Q^{(c)}(\omega) = \left| p(\omega) \right|^2,
\end{equation}

\noindent with $p(\omega)$ evaluated conformly to Eq. (\ref{eqn:additional}).



\end{itemize}
\end{textblock*}



\end{frame}






\begin{frame}
\frametitle{Quantum theory of raidation by nonstationary systems 
	with application to high-order hamronic generation
        (Phys. Rev. A 101, 013410 (2020))}
	\framesubtitle{Mathematical details}
\small

\begin{textblock*}{400pt}(30pt,60pt)
\footnotesize
\begin{itemize}
\item Integration is done by using the following Simpson rule,
$m$ being the number of subintervals
with $h=(x_{m}-x_0)/m$, $x_i = x_0 + ih$,
\begin{align}
\begin{split}
 \int_{x_0}^{x_m} f(x)dx & \sim \frac{1}{3} h \sum_{i=1}^{m/2}
  \left[ f(x_{2i-2})  + 4f(x_{2i-1}) + f(x_{2i})\right]\;\;\mbox{ (1/3 rule) }
\end{split} \label{eqn:simps18} 
%\begin{split}
% \int_{x_0}^{x_m} f(x)dx & \sim \frac{3}{8} h \sum_{i=1}^{m/3}
%  \left[ f(x_{3i-3})  + 3f(x_{3i-2}) +    \right.
%  \left. 3f(x_{3i-1}) + f(x_{3i})\right]\;\;\mbox{ (3/8 rule) },
%\end{split} \label{eqn:simps38} 
\end{align}

or the trapezoidal rule, given by
\begin{align}
\begin{split}
 \int_{x_0}^{x_m} f(x)dx 
   & \sim \sum_{i=1}^{m} \frac{ f(x_{i-1})  + f(x_{i})}{2}\Delta x_i
\end{split} \label{eqn:trapez} 
\end{align}

\item Numerical differentiation may be done by using the formulas

\begin{align}
\begin{split}
f^{\prime}_l (x_i) = \frac{f(x_{i+1}) -  f(x_i)}{x_{i+1} -x_i}, & \hspace{20pt}
f^{\prime}_r (x_i) = \frac{f(x_{i}) -  f(x_{i-1})}{x_{i} - x_{i-1}}
\end{split} \label{eqn:diff}   \\
f^{\prime} (x_i) = & \frac{f^{\prime}_l (x_i) + f^{\prime}_r (x_i)}{2}
%\begin{split}
% \int_{x_0}^{x_m} f(x)dx & \sim \frac{3}{8} h \sum_{i=1}^{m/3}
%  \left[ f(x_{3i-3})  + 3f(x_{3i-2}) +    \right.
%  \left. 3f(x_{3i-1}) + f(x_{3i})\right]\;\;\mbox{ (3/8 rule) },
%\end{split} \label{eqn:simps38} 
\end{align}

  
\end{itemize}

\end{textblock*}


\end{frame}







\begin{frame}
\frametitle{Quantum theory of raidation by nonstationary systems 
	with application to high-order hamronic generation
        (Phys. Rev. A 101, 013410 (2020))}
	\framesubtitle{Time dependence - $T=3\tau$}
\small
\begin{textblock*}{500pt}(5pt,60pt)
%\centering
%\rotatebox{90}{
 \begin{minipage}{0.0\linewidth}
 \begin{figure}[tp]
 \includegraphics[scale=.20, trim={6 100 10 133},clip]{images/114_f01_3_density_comp.png} \\
 \includegraphics[scale=.20, trim={6 100 10 130},clip]{images/114_f01_3_hhg_comp.png} \\
% \includegraphics[scale=.20, trim={6 130 10 130},clip]{images/114_f003_8_03_hhg.png} 
%%\caption{Eigenvalues.}
%\label{fig:pemd_0114}
 \end{figure}

\end{minipage}
\end{textblock*}


\begin{textblock*}{500pt}(130pt,60pt)
%\centering
%\rotatebox{90}{
 \begin{minipage}{0.0\linewidth}
 \begin{figure}[tp]
 \includegraphics[scale=.20, trim={6 100 10 133},clip]{images/114_f01_3_values_100.png} \\
 \includegraphics[scale=.20, trim={6 100 10 130},clip]{images/114_f01_3_values_150.png} \\
% \includegraphics[scale=.20, trim={30 130 10 130},clip]{images/114_f003_12_07_hhg.png} 
%%\caption{Eigenvalues.}
%\label{fig:pemd_0114}
 \end{figure}

\end{minipage}
\end{textblock*}


\begin{textblock*}{500pt}(255pt,60pt)
%\centering
%\rotatebox{90}{
 \begin{minipage}{0.0\linewidth}
 \begin{figure}[tp]
 \includegraphics[scale=.20, trim={6 100 10 133},clip]{images/114_f01_3_values_200.png} \\
 \includegraphics[scale=.20, trim={6 100 10 130},clip]{images/114_f01_3_values_209.png} \\
% \includegraphics[scale=.20, trim={30 130 10 130},clip]{images/114_f003_12_07_hhg.png} 
%%\caption{Eigenvalues.}
%\label{fig:pemd_0114}
 \end{figure}

\end{minipage}
\end{textblock*}


\begin{textblock*}{500pt}(390pt,165pt)
%\centering
\rotatebox{90}{
 \begin{minipage}{0.0\linewidth}
 \begin{figure}[tp]
 \includegraphics[scale=.20, trim={6 100 10 120},clip]{images/114_f01_3_error.png } \\
%%\caption{Eigenvalues.}
%\label{fig:pemd_0114}
 \end{figure}
\end{minipage}
}
\end{textblock*}




\begin{textblock*}{\linewidth}(20pt,170pt)
\rotatebox{0}{
\begin{minipage}{\linewidth}
{
\tiny 
\centering
\input{tables/114_f01_3t_error_2.tab}
}
\end{minipage}
}
\end{textblock*}


\begin{textblock*}{400pt}(30pt,205pt)
\scriptsize

\begin{align}
    \sigma_{ij}= \left|\frac{\left(   Q^{(c)}(\omega)[i] 
       - Q^{(c)}(\omega)[j]  \right)}{Q^{(c)}(\omega)[j]}\right|, \hspace{20pt}
    Q^{(c)}(\omega) & = \left|\Bigg[ e^{i\omega t} x(t)\Bigg]_{-T/2}^{T/2}
               -i\omega \int_{-T/2}^{T/2} e^{i\omega t}x(t)dt\right|^2.
                \nonumber
\end{align}

\end{textblock*}

\end{frame}






\begin{frame}
\frametitle{Quantum theory of raidation by nonstationary systems 
	with application to high-order hamronic generation
        (Phys. Rev. A 101, 013410 (2020))}
	\framesubtitle{Time dependence - $T=5\tau$}
\small
\begin{textblock*}{500pt}(5pt,60pt)
%\centering
%\rotatebox{90}{
 \begin{minipage}{0.0\linewidth}
 \begin{figure}[tp]
 \includegraphics[scale=.20, trim={6 100 10 133},clip]{images/114_f01_5_density_comp.png} \\
 \includegraphics[scale=.20, trim={6 100 10 130},clip]{images/114_f01_5_hhg_comp.png} \\
% \includegraphics[scale=.20, trim={6 130 10 130},clip]{images/114_f003_8_03_hhg.png} 
%%\caption{Eigenvalues.}
%\label{fig:pemd_0114}
 \end{figure}

\end{minipage}
\end{textblock*}


\begin{textblock*}{500pt}(130pt,60pt)
%\centering
%\rotatebox{90}{
 \begin{minipage}{0.0\linewidth}
 \begin{figure}[tp]
 \includegraphics[scale=.20, trim={6 100 10 133},clip]{images/114_f01_5_values_100.png} \\
 \includegraphics[scale=.20, trim={6 100 10 130},clip]{images/114_f01_5_values_150.png} \\
% \includegraphics[scale=.20, trim={30 130 10 130},clip]{images/114_f003_12_07_hhg.png} 
%%\caption{Eigenvalues.}
%\label{fig:pemd_0114}
 \end{figure}

\end{minipage}
\end{textblock*}


\begin{textblock*}{500pt}(255pt,60pt)
%\centering
%\rotatebox{90}{
 \begin{minipage}{0.0\linewidth}
 \begin{figure}[tp]
 \includegraphics[scale=.20, trim={6 100 10 133},clip]{images/114_f01_5_values_200.png} \\
 \includegraphics[scale=.20, trim={6 100 10 130},clip]{images/114_f01_5_values_209.png} \\
% \includegraphics[scale=.20, trim={30 130 10 130},clip]{images/114_f003_12_07_hhg.png} 
%%\caption{Eigenvalues.}
%\label{fig:pemd_0114}
 \end{figure}

\end{minipage}
\end{textblock*}


\begin{textblock*}{500pt}(390pt,165pt)
%\centering
\rotatebox{90}{
 \begin{minipage}{0.0\linewidth}
 \begin{figure}[tp]
 \includegraphics[scale=.20, trim={6 100 10 132},clip]{images/114_f01_5_error.png } \\
%%\caption{Eigenvalues.}
%\label{fig:pemd_0114}
 \end{figure}
\end{minipage}
}
\end{textblock*}




\begin{textblock*}{\linewidth}(20pt,170pt)
\rotatebox{0}{
\begin{minipage}{\linewidth}
{
\tiny 
\centering
\input{tables/114_f01_5t_error_2.tab}
}
\end{minipage}
}
\end{textblock*}


\begin{textblock*}{400pt}(30pt,205pt)
\scriptsize

\begin{align}
    \sigma_{ij} = \left|\frac{\left(   Q^{(c)}(\omega)[i] 
       - Q^{(c)}(\omega)[j]  \right)}{Q^{(c)}(\omega)[j]}\right|, \hspace{20pt}
    Q^{(c)}(\omega) & = \left|\Bigg[ e^{i\omega t} x(t)\Bigg]_{-T/2}^{T/2}
               -i\omega \int_{-T/2}^{T/2} e^{i\omega t}x(t)dt\right|^2.
                \nonumber
\end{align}

\end{textblock*}

\end{frame}





\begin{frame}
\frametitle{Quantum theory of raidation by nonstationary systems 
	with application to high-order hamronic generation
        (Phys. Rev. A 101, 013410 (2020))}
	\framesubtitle{Time dependence - $T=8\tau$}
\small
\begin{textblock*}{500pt}(5pt,60pt)
%\centering
%\rotatebox{90}{
 \begin{minipage}{0.0\linewidth}
 \begin{figure}[tp]
 \includegraphics[scale=.20, trim={6 100 10 133},clip]{images/114_f01_8_density_comp.png} \\
 \includegraphics[scale=.20, trim={6 100 10 130},clip]{images/114_f01_8_hhg_comp.png} \\
% \includegraphics[scale=.20, trim={6 130 10 130},clip]{images/114_f003_8_03_hhg.png} 
%%\caption{Eigenvalues.}
%\label{fig:pemd_0114}
 \end{figure}

\end{minipage}
\end{textblock*}


\begin{textblock*}{500pt}(130pt,60pt)
%\centering
%\rotatebox{90}{
 \begin{minipage}{0.0\linewidth}
 \begin{figure}[tp]
 \includegraphics[scale=.20, trim={6 100 10 133},clip]{images/114_f01_8_values_100.png} \\
 \includegraphics[scale=.20, trim={6 100 10 130},clip]{images/114_f01_8_values_150.png} \\
% \includegraphics[scale=.20, trim={30 130 10 130},clip]{images/114_f003_12_07_hhg.png} 
%%\caption{Eigenvalues.}
%\label{fig:pemd_0114}
 \end{figure}

\end{minipage}
\end{textblock*}


\begin{textblock*}{500pt}(255pt,60pt)
%\centering
%\rotatebox{90}{
 \begin{minipage}{0.0\linewidth}
 \begin{figure}[tp]
 \includegraphics[scale=.20, trim={6 100 10 133},clip]{images/114_f01_8_values_200.png} \\
 \includegraphics[scale=.20, trim={6 100 10 130},clip]{images/114_f01_8_values_209.png} \\
% \includegraphics[scale=.20, trim={30 130 10 130},clip]{images/114_f003_12_07_hhg.png} 
%%\caption{Eigenvalues.}
%\label{fig:pemd_0114}
 \end{figure}

\end{minipage}
\end{textblock*}


\begin{textblock*}{500pt}(390pt,165pt)
%\centering
\rotatebox{90}{
 \begin{minipage}{0.0\linewidth}
 \begin{figure}[tp]
 \includegraphics[scale=.20, trim={6 100 10 120},clip]{images/114_f01_8_error.png } \\
%%\caption{Eigenvalues.}
%\label{fig:pemd_0114}
 \end{figure}
\end{minipage}
}
\end{textblock*}




\begin{textblock*}{\linewidth}(20pt,170pt)
\rotatebox{0}{
\begin{minipage}{\linewidth}
{
\tiny 
\centering
\input{tables/114_f01_8t_error_2.tab}
}
\end{minipage}
}
\end{textblock*}


\begin{textblock*}{400pt}(30pt,205pt)
\scriptsize

\begin{align}
    \sigma_{ij} = \left|\frac{\left(   Q^{(c)}(\omega)[i] 
       - Q^{(c)}(\omega) [j] \right)}{Q^{(c)}(\omega)[j]}\right|, \hspace{20pt}
    Q^{(c)}(\omega) & = \left|\Bigg[ e^{i\omega t} x(t)\Bigg]_{-T/2}^{T/2}
               -i\omega \int_{-T/2}^{T/2} e^{i\omega t}x(t)dt\right|^2.
                \nonumber
\end{align}

\end{textblock*}

\end{frame}



\begin{frame}
\frametitle{Quantum theory of raidation by nonstationary systems 
	with application to high-order hamronic generation
        (Phys. Rev. A 101, 013410 (2020))}
	\framesubtitle{Time dependence reference table}
\small

\begin{textblock*}{400pt}(30pt,60pt)
\footnotesize
\begin{itemize}
\item Sufficient time step for converged result is 
$\Delta t \leqslant 0.08\,\rm a.u.$
\item In the following, we will maintain $\Delta t$ set to 
around $0.04\,\rm a.u.$
\end{itemize}

\end{textblock*}

{
\tiny 
\input{tables/ref_dt.tab}
}


\begin{textblock*}{400pt}(30pt,205pt)
\scriptsize

\begin{align}
    \tau = n_{oc}\frac{2\pi}{\omega}, \hspace{20pt}
    T = n \tau, \hspace{20pt}
    \Delta t  = \frac{T}{{\rm nsteps}}, \hspace{20pt}
    n_{oc}  = 5,
                \nonumber
\end{align}

\end{textblock*}

\end{frame}




\begin{frame}
\frametitle{Quantum theory of raidation by nonstationary systems 
	with application to high-order hamronic generation
        (Phys. Rev. A 101, 013410 (2020))}
	\framesubtitle{Time dependence - $\rm{nsteps} = 100000$}
\small
\begin{textblock*}{500pt}(5pt,60pt)
%\centering
%\rotatebox{90}{
 \begin{minipage}{0.0\linewidth}
 \begin{figure}[tp]
 \includegraphics[scale=.20, trim={6 100 10 133},clip]{images/114_f01_T_density_comp.png} \\
 \includegraphics[scale=.20, trim={6 100 10 130},clip]{images/114_f01_T_hhg_comp.png} \\
% \includegraphics[scale=.20, trim={6 130 10 130},clip]{images/114_f003_8_03_hhg.png} 
%%\caption{Eigenvalues.}
%\label{fig:pemd_0114}
 \end{figure}

\end{minipage}
\end{textblock*}


\begin{textblock*}{500pt}(130pt,60pt)
%\centering
%\rotatebox{90}{
 \begin{minipage}{0.0\linewidth}
 \begin{figure}[tp]
 \includegraphics[scale=.20, trim={6 100 10 133},clip]{images/114_f01_T_values_100.png} \\
 \includegraphics[scale=.20, trim={6 100 10 130},clip]{images/114_f01_T_values_150.png} \\
% \includegraphics[scale=.20, trim={30 130 10 130},clip]{images/114_f003_12_07_hhg.png} 
%%\caption{Eigenvalues.}
%\label{fig:pemd_0114}
 \end{figure}

\end{minipage}
\end{textblock*}


\begin{textblock*}{500pt}(255pt,60pt)
%\centering
%\rotatebox{90}{
 \begin{minipage}{0.0\linewidth}
 \begin{figure}[tp]
 \includegraphics[scale=.20, trim={6 100 10 133},clip]{images/114_f01_T_values_200.png} \\
 \includegraphics[scale=.20, trim={6 100 10 130},clip]{images/114_f01_T_values_209.png} \\
% \includegraphics[scale=.20, trim={30 130 10 130},clip]{images/114_f003_12_07_hhg.png} 
%%\caption{Eigenvalues.}
%\label{fig:pemd_0114}
 \end{figure}

\end{minipage}
\end{textblock*}


%\begin{textblock*}{500pt}(390pt,165pt)
%%\centering
%\rotatebox{90}{
% \begin{minipage}{0.0\linewidth}
% \begin{figure}[tp]
% \includegraphics[scale=.20, trim={6 100 10 120},clip]{images/114_f01_8_error.png } \\
%%%\caption{Eigenvalues.}
%%\label{fig:pemd_0114}
% \end{figure}
%\end{minipage}
%}
%\end{textblock*}




\begin{textblock*}{\linewidth}(20pt,170pt)
\rotatebox{0}{
\begin{minipage}{\linewidth}
{
\tiny 
\centering
\input{tables/114_f01_T_100000.tab}
}
\end{minipage}
}
\end{textblock*}


\begin{textblock*}{400pt}(30pt,205pt)
\scriptsize

\begin{align}
%   \sigma_{ij} = \left|\frac{\left(   Q^{(c)}(\omega)[i] 
%      - Q^{(c)}(\omega) [j] \right)}{Q^{(c)}(\omega)[j]}\right|, \hspace{20pt}
    Q^{(c)}(\omega) & = \left|\Bigg[ e^{i\omega t} x(t)\Bigg]_{-T/2}^{T/2}
               -i\omega \int_{-T/2}^{T/2} e^{i\omega t}x(t)dt\right|^2.
                \nonumber
\end{align}

\end{textblock*}

\end{frame}










\begin{frame}
\frametitle{Quantum theory of raidation by nonstationary systems 
	with application to high-order hamronic generation
        (Phys. Rev. A 101, 013410 (2020))}
	\framesubtitle{Time dependence - $\Delta t= 0.04\,\rm a.u.$}
\small
\begin{textblock*}{500pt}(5pt,60pt)
%\centering
%\rotatebox{90}{
 \begin{minipage}{0.0\linewidth}
 \begin{figure}[tp]
 \includegraphics[scale=.20, trim={6 100 10 133},clip]{images/114_f01_T_density_comp_2.png} \\
 \includegraphics[scale=.20, trim={6 100 10 130},clip]{images/114_f01_T_hhg_comp_2.png} \\
% \includegraphics[scale=.20, trim={6 130 10 130},clip]{images/114_f003_8_03_hhg.png} 
%%\caption{Eigenvalues.}
%\label{fig:pemd_0114}
 \end{figure}

\end{minipage}
\end{textblock*}


\begin{textblock*}{500pt}(130pt,60pt)
%\centering
%\rotatebox{90}{
 \begin{minipage}{0.0\linewidth}
 \begin{figure}[tp]
 \includegraphics[scale=.20, trim={6 100 10 133},clip]{images/114_f01_T_values_dt_100.png} \\
 \includegraphics[scale=.20, trim={6 100 10 130},clip]{images/114_f01_T_values_dt_150.png} \\
% \includegraphics[scale=.20, trim={30 130 10 130},clip]{images/114_f003_12_07_hhg.png} 
%%\caption{Eigenvalues.}
%\label{fig:pemd_0114}
 \end{figure}

\end{minipage}
\end{textblock*}


\begin{textblock*}{500pt}(255pt,60pt)
%\centering
%\rotatebox{90}{
 \begin{minipage}{0.0\linewidth}
 \begin{figure}[tp]
 \includegraphics[scale=.20, trim={6 100 10 133},clip]{images/114_f01_T_values_dt_200.png} \\
 \includegraphics[scale=.20, trim={6 100 10 130},clip]{images/114_f01_T_values_dt_209.png} \\
% \includegraphics[scale=.20, trim={30 130 10 130},clip]{images/114_f003_12_07_hhg.png} 
%%\caption{Eigenvalues.}
%\label{fig:pemd_0114}
 \end{figure}

\end{minipage}
\end{textblock*}


%\begin{textblock*}{500pt}(390pt,165pt)
%%\centering
%\rotatebox{90}{
% \begin{minipage}{0.0\linewidth}
% \begin{figure}[tp]
% \includegraphics[scale=.20, trim={6 100 10 120},clip]{images/114_f01_8_error.png } \\
%%%\caption{Eigenvalues.}
%%\label{fig:pemd_0114}
% \end{figure}
%\end{minipage}
%}
%\end{textblock*}




\begin{textblock*}{\linewidth}(20pt,170pt)
\rotatebox{0}{
\begin{minipage}{\linewidth}
{
\tiny 
\centering
\input{tables/114_f01_T_dt004.tab}
}
\end{minipage}
}
\end{textblock*}


\begin{textblock*}{400pt}(30pt,205pt)
\scriptsize

\begin{align}
%   \sigma_{ij} = \left|\frac{\left(   Q^{(c)}(\omega)[i] 
%      - Q^{(c)}(\omega) [j] \right)}{Q^{(c)}(\omega)[j]}\right|, \hspace{20pt}
    Q^{(c)}(\omega) & = \left|\Bigg[ e^{i\omega t} x(t)\Bigg]_{-T/2}^{T/2}
               -i\omega \int_{-T/2}^{T/2} e^{i\omega t}x(t)dt\right|^2.
                \nonumber
\end{align}

\end{textblock*}

\end{frame}



\begin{frame}
\frametitle{Quantum theory of raidation by nonstationary systems 
	with application to high-order hamronic generation
        (Phys. Rev. A 101, 013410 (2020))}
	\framesubtitle{Convergence parameters table for 
             the homogeneous equation}
\small

\begin{textblock*}{400pt}(30pt,60pt)
\footnotesize
\begin{itemize}
\item The following table shows parameters known to return converged results 
for Eq. (\ref{eqn:step_1d}), which is recalled below
%\item \textbf{For $\omega=0.057\,\rm a.u.$ and $F_0 = 0.3\,\rm a.u.$,
%PEMD and HHG spectra currently do not converge even 
%when density probability does}
\end{itemize}

\end{textblock*}

\vspace{20pt}

{
\tiny 
\input{tables/conv_table_057.tab}
\input{tables/conv_table_114.tab}
}


\begin{textblock*}{400pt}(30pt,205pt)
\scriptsize

\begin{equation}
	i\frac{\partial}{\partial t}\psi(x,t)
	   = \left( \hat{H}_a + \hat{V}_{\rm ext}(t)\right) 
	   \psi(x,t),   \hspace{20pt}\psi(x,0)=\psi_0(x), \nonumber
\end{equation}
\begin{equation}
  \hat{H}_a = -\frac{1}{2}\frac{d^2}{dx^2} - \varkappa\delta(x), \hspace{20pt}
  \hat{V}_{\rm ext} = F(t)x,   \nonumber
\end{equation}

\end{textblock*}


\begin{textblock*}{100pt}(30pt,140pt)
\scriptsize
\textbf{Legend} \\
$\rm (nbas/nelm) - xmax $


\end{textblock*}

\end{frame}



\begin{frame}
\frametitle{Quantum theory of raidation by nonstationary systems 
	with application to high-order hamronic generation
        (Phys. Rev. A 101, 013410 (2020))}
	\framesubtitle{Structure of the code, performance 
         and memory management}
\small

\begin{textblock*}{400pt}(30pt,60pt)
\footnotesize
\begin{enumerate}
\item Diagonalization
\begin{itemize}
\footnotesize
\item[$\bullet$] most time consuming part for bigger matrixes
\item[$\bullet$] is there a way to multithread lapack diagonalization? 
\end{itemize}
\item Propagation
\begin{itemize}
\footnotesize
\item[$\bullet$] can be arbitarily long 
(convergence requires $\Delta t \leqslant 0.08\,\rm a.u.$)
\item[$\bullet$] order 4 scheme required
\item[$\bullet$] multithreading would also help here
\end{itemize}
\item Projection 
\begin{itemize}
\footnotesize
\item[$\bullet$] very short (matter of minutes)
\item[$\bullet$] potentially very memory consuming
\end{itemize}
\end{enumerate}

\end{textblock*}

\end{frame}


\begin{frame}
\frametitle{Quantum theory of raidation by nonstationary systems 
	with application to high-order hamronic generation
        (Phys. Rev. A 101, 013410 (2020))}
	\framesubtitle{What is being addressed right now}
\small

\begin{textblock*}{400pt}(30pt,60pt)
\footnotesize
\begin{itemize}
\item Fixing the code for performance
\begin{itemize}
\footnotesize
\item[$\bullet$] At the current state, for $\omega=0.057\,\rm a.u.$
and $F_0=0.3\,$, PEMD and HHG-C spectra are wrong even if 
the probability density seems right
\item[$\bullet$] It currently make more sense to totally separate 
propagation and observables computation programs
\end{itemize}
\item Inhomogeneous term should be addressed
\begin{enumerate}
\footnotesize
\item[$\bullet$] Find a proper Runge Kutta case in the litterature
\item[$\bullet$] Use Runge Kutta from Numerical Recipes as a reference
\item[$\bullet$] Write own code and compare 
\end{enumerate}

\end{itemize}

\end{textblock*}

\end{frame}







%\begin{frame}
%\frametitle{Unveiling the Inhomogeneous Nature of 
%	Strong Field Ionization in Extended Systems 
%        (Phys. Rev. A 101, 013410 (2020))}
%	\framesubtitle{The inhomogeneous equation}
%\small
%
%\begin{textblock*}{400pt}(30pt,60pt)
%\footnotesize
%\begin{itemize}
%\item We know want by the means of RK4 method, to deal with Eq. (\ref{eqn:main_1d}), i.e.
%
%\begin{equation}
%	i\frac{\partial}{\partial t}\phi_\omega(x,t)
%	   = \left( \hat{H}_a + \hat{V}_{\rm ext}(t)\right) 
%	   \phi_\omega (x,t)
% 	   + e^{i\omega t} {\hat{p}}\,
%	   \psi(x,t), \nonumber
%\end{equation}
%\begin{equation}
%  \phi_\omega (x,0) = \frac{-i\varkappa^{3/2}}{\omega}
%{\rm sgn}(x)\left(e^{-\varkappa|x|} - e^{-\varkappa(\omega)|x|}\right), 
%\hspace{20pt}  \psi(x,0) = \varkappa^{1/2}e^{-\varkappa|x|}, 
%\nonumber
%\end{equation}
%\item To this end, we intend to do the following
%  \begin{enumerate}
%    \footnotesize
%    \item Successfully implement Runge Kutta from 
%       "Numerical Recipes in Fortran" to our problem
%    \item Solve Eq. (\ref{eqn:main_1d}) with any scheme from 
%   Kidd \textit{et al.}, Phys, Rev E 96, 063307 (2017)
%    \item Fix our own implementation of RK4
%  \end{enumerate}
%\end{itemize}
%
%\end{textblock*}
%
%\end{frame}
